%! Author = lili
%! Date = 6/16/26

\section*{Information Extraction from Structured, Semi-Structured, and Unstructured Data}

\section{Motivation and Background}
\defin{kga}


\begin{center}
    ``\textit{Learning is associated with increased degrees of true belief, although these beliefs need not
    necessarily be justified. Knowledge acquisition additionally requires that the true belief is
    justified, or has been formed for the right reason.}''
\end{center}
- Hössjer, Díaz-Pachón, and Rao 2022

\begin{mytable}
    \begin{tabular}{p{0.45\textwidth} p{0.45\textwidth}}

        \textbf{KG Reasoning} & \textbf{KG Acquisition} \\

        Assumes KG already exists &
        Starts from raw data (text, tables, etc.) \\

        Focus: Inferring new facts from existing ones &
        Focus: Building the KG from scratch \\

        Methods: Logical rules, embeddings, neural models &
        Methods: NER, NED, NEL, LLM-based extraction \\

    \end{tabular}
    \caption{Comparison of KG Reasoning and KG Acquisition.}
    \label{tab:kg_reasoning_vs_acquisition}
\end{mytable}

\[
    \text{Data} \overset{\text{Information Extraction}}{\Rightarrow} \text{KG} \overset{\text{Reasoning}}{\Rightarrow} \text{New Knowledge}
\]

\begin{center}
    ``\textit{Knowledge extraction can be performed from structured,
        semi-structured, and unstructured data sources.}''
\end{center}
- Zhu et al. 2024)

\begin{mytable}
    \begin{tabular}{p{0.3\textwidth} p{0.6\textwidth}}
        \toprule
        \textbf{Data Type} & \textbf{Description} \\
        \midrule
        Structured & Relational databases, tables with predefined schemas \\

        Semi-Structured & XML, JSON, HTML, Wikipedia infoboxes, documents with some structural markup \\

        Unstructured & Raw text from news articles, social media posts, scientific papers, web pages \\
        \bottomrule
    \end{tabular}
    \caption{The three types of data sources for knowledge graph construction.}
    \label{tab:data_types}
\end{mytable}

\begin{challenge}
    ``\textit{The majority of previous studies have focused on unstructured text, as this represents the
    most abundant and challenging data source for KG construction.}''
\end{challenge}
- Al-Moslmi et al. 2020

\paragraph{Named entity extraction} has three core tasks:
\defin{ner}
\defin{ned}
\defin{nel}

\begin{bsp}
    \defiref{ner} \\
    find every segment in a text that
    mentions a named entity
\end{bsp}

\begin{bsp}
    \defiref{ned} \\
    determine which named entity a
    mention refers to
\end{bsp}

\begin{bsp}
    \defiref{nel} \\
    provide a standard \gls{iri} for each
    disambiguated entity
\end{bsp}

\subsection*{Structured Data: Direct Mapping to KGs}
\begin{center}
    ``\textit{Representing NL text in computer-processable form, KGs are a natural way: entities
    can be represented as nodes in KGs and relations as edges.}''
\end{center}
- Al-Moslmi et al. 2020

\begin{bsp}
\begin{mytable}
    \begin{tabular}{cccc}
        \toprule
        \textbf{Event} & \textbf{Date} & \textbf{Attribute-1} & \textbf{Attribute-2} \\
        \midrule
        1 & 14 Mar & severe & 38.5 \\
        2 & 18 Jun & mild & 37.5 \\
        3 & 23 Jul & mild & 36.9 \\
        \bottomrule
    \end{tabular}
    \caption{Time-Oriented Database Example.}
    \label{tab:event_attributes}
\end{mytable}
    As Triple:
    \[
        <Event1> <hasSeverity> <severe>
        <Event1> <hasTemp> <38.5>
    \]
\end{bsp}

``The defining
characteristic of a time-oriented database is that
sequential values for each attribute may be
recorded. The time intervals between events
may differ.''

The advantages of direct mapping is that there is no ambiguity and predicates are defined by clear schemas.

\subsection*{Semi-Structured Data}
\begin{bsp}
    \textbf{Wikipedia Infobox} \\
    Wikipedia infoboxes contain attribute-value pairs that can be directly mapped to KG
    triples.
\end{bsp}
